% ============================================================================
% PART 08 - FUNCTIONAL CLASSIFICATION
% ============================================================================

\labelsec{functional_classification}

This chapter demonstrates that the three SHT units encode sufficient information to classify DNA functional state with high accuracy.

\section{Classification Hypothesis}

\begin{keyresult}
\textbf{Hypothesis:} The geometric signature $[v_0,\,\eta_r,\,\Omega]$ distinguishes DNA functional class without any sequence information.
\end{keyresult}

This is non-trivial: conventional descriptors like GC content or length differ between classes but are not geometric. SHT bypasses sequence and directly characterizes 3D structure.

\section{Training Dataset}

\begin{table}[h]
    \centering
    \caption{Labeled Structures for Classification ($N=33$)}\labeltab{class_labels}
    \begin{tabular}{llll}
        \toprule
        \textbf{Class} & $N$ & \textbf{Structures} & \textbf{Known signature} \\
        \midrule
        Nucleosome      & 2  & 1AOI, 1F66           & High Vij, high Mamton \\
        Protein-bound   & 11 & Various TF complexes  & Moderate Vij, variable Mamton \\
        Free duplex     & 16 & Canonical B-DNA       & Low Vij, low Mamton \\
        Non-canonical   & 4  & G4, damage, junction  & Extreme Shambhian \\
        \bottomrule
        \multicolumn{2}{l}{\textbf{Total}} & 33 & \\
        \bottomrule
    \end{tabular}
\end{table}

Labels were assigned from PDB metadata \textit{before} examining geometric values. The remaining 167 structures are unlabeled and serve as prediction targets.

\section{Model Architecture}

A Random Forest Classifier was chosen for:
\begin{itemize}
    \item Robustness to scale differences between units (no normalization needed)
    \item Native feature importance via Gini impurity
    \item Non-linear decision boundaries without manual kernel choice
    \item Calibrated probability outputs for confidence scoring
\end{itemize}

\textbf{Hyperparameters:} 200 trees, max depth 10, min samples per leaf 1, Gini criterion. No grid search was performed (small $N = 33$); default scikit-learn parameters were used.

\section{Performance Results}

\subsection{Cross-Validation}

\begin{keyresult}
    \textbf{5-fold stratified cross-validation accuracy:} $\mathbf{91.0\% \pm 11.7\%}$
    
    Tested on $N = 33$ labeled structures. Chance performance for 4 classes = 25\%.
\end{keyresult}

The high standard deviation ($\pm 11.7\%$) reflects the small sample size ($N = 33$, $\sim$6--7 per fold) and the class imbalance (only 2 nucleosome structures). With a larger dataset, variance would decrease.

\subsection{Per-Class Performance (Estimated)}

\begin{table}[h]
    \centering
    \caption{Estimated Per-Class Precision and Recall}\labeltab{per_class}
    \begin{tabular}{lllll}
        \toprule
        \textbf{Class} & $N$ & \textbf{Precision} & \textbf{Recall} & \textbf{Comments} \\
        \midrule
        Nucleosome    & 2  & $\sim$1.00 & $\sim$1.00 & Extreme outliers, easily separated \\
        Protein-bound & 11 & $\sim$0.85 & $\sim$0.85 & Some overlap with free duplex \\
        Free duplex   & 16 & $\sim$0.94 & $\sim$0.94 & Largest class; cleanest boundary \\
        Non-canonical & 4  & $\sim$0.75 & $\sim$0.75 & Heterogeneous class \\
        \bottomrule
    \end{tabular}
\end{table}

\section{Feature Importance}

The trained forest assigns the following Gini-importance to each feature:

\begin{table}[h]
    \centering
    \caption{SHT Feature Importance (Random Forest)}\labeltab{feature_importance}
    \begin{tabular}{lll}
        \toprule
        \textbf{Feature} & \textbf{Importance} & \textbf{Interpretation} \\
        \midrule
        Mamton ($\Omega$) & \textbf{41\%} & Most discriminative; extreme for nucleosomes \\
        Vij ($v_0$)       & 33\%          & Second; separates bent from straight DNA \\
        Shambhian ($\eta_r$) & 26\%       & Least; useful for non-canonical class \\
        \bottomrule
    \end{tabular}
\end{table}

Mamton is the most important feature despite being a derived quantity (3rd-order moment). Its dominance reflects the extreme dynamic range: $\Omega$ spans 3.7 orders of magnitude across the dataset, providing strong class separation.

\begin{figure}[h]
    \centering
    \includegraphics[width=0.65\textwidth]{feature_importance}
    \caption{Random Forest feature importance (Gini impurity). Mamton is the most discriminative unit (41\%), followed by Shambhian (33\%) and Vij (26\%).}\labelfig{feature_importance}
\end{figure}

\section{Blind Predictions}

\subsection{Case Study: Structure 1F66}

Before looking up PDB metadata, 1F66 was presented to the classifier as unlabeled. The feature vector was $[v_0 = 0.297,\;\eta_r = 0.828,\;\Omega = 26{,}441]$.

\begin{keyresult}
    \textbf{Prediction:} Nucleosome (probability $> 0.99$) \\
    \textbf{Confirmation:} H2A.Z-containing nucleosome core particle \citep{Suto2000} \\
    \textbf{Verdict:} Correct --- blind prediction confirmed by PDB record
\end{keyresult}

The z-scores of 1F66 make this prediction unsurprising in retrospect: both Vij and Mamton are $+4$--$+5\sigma$ above the dataset mean. But crucially, this structure was \textbf{not in the training set} and was classified on geometry alone.

\subsection{Predictions on Unlabeled Structures}

The remaining 167 unlabeled structures were classified. Key predictions:

\begin{itemize}
    \item \textbf{Free duplex} (predicted): 89 structures --- standard B-DNA, consistent with crystallography conditions
    \item \textbf{Protein-bound} (predicted): 58 structures --- various complexes, likely overfit slightly given small training set
    \item \textbf{Non-canonical} (predicted): 18 structures --- geometric anomalies worth experimental follow-up
    \item \textbf{Nucleosome candidate} (predicted, extreme $\Omega$): 2 structures besides 1AOI/1F66 --- merit curated PDB annotation check
\end{itemize}

These predictions are exploratory. Experimental verification and larger labeled datasets are needed before production use.

\section{Decision Boundaries and Geometric Interpretation}

The classifier's first branching level (root node of most trees) uses a Mamton threshold at $\Omega \approx 5{,}000$: structures above this are almost exclusively protein-wrapped or nucleosomal. The second level uses Vij to distinguish protein-bound from nucleosomal. Shambhian is most active in the final level to separate non-canonical from free-duplex.

\begin{figure}[h]
    \centering
    \includegraphics[width=0.85\textwidth]{functional_classification_space}
    \caption{Functional classification in SHT geometric space. Points are colored by predicted class. Nucleosome structures (1AOI, 1F66) are visually isolated in the extreme Vij--Mamton corner. Free duplex structures cluster at low Mamton. The PCA projection collapses three dimensions to two for visualization.}\labelfig{functional_class_space}
\end{figure}

This three-level hierarchy maps naturally onto biological intuition: \textit{how rough is it?} $\rightarrow$ \textit{how bent is it?} $\rightarrow$ \textit{how distorted is the cross-section?}

